\chapter{Manual de usuario}
\label{app:manual_usuario}

\section{Acceso al sistema}
\label{sec:acceso}


La versión del generador descrita en este manual puede utilizarse ejecutando la
aplicación en un entorno local, siguiendo los pasos indicados en la
Sección~\ref{sec:verificacion-instalacion} del manual de instalación y
despliegue.

UVLHub dispone de una instancia pública en \url{https://www.uvlhub.io/}.
No obstante, esta instancia puede evolucionar de forma independiente al código
descrito en esta memoria. Por este motivo, la ejecución local mediante Docker
constituye la forma de acceso reproducible a la versión evaluada en el presente
trabajo.

La generación de modelos no requiere que el usuario disponga de una cuenta ni
que realice un proceso de autenticación previo.

La pantalla inicial de UVLHub se muestra en la Figura~\ref{fig:inicio-uvlhub}.

\begin{figure}[H]
    \centering
    \makebox[\textwidth][c]{%
        \includegraphics[width=1.2\textwidth]{inicio_uvlhub.png}
    }
    \caption{Pantalla inicial de UVLHub. Captura realizada por el autor.}
    \label{fig:inicio-uvlhub}
\end{figure}

\section{Funcionalidad principal}
\label{sec:funcionalidad_principal}


La funcionalidad principal del sistema consiste en la generación automática de
modelos de características mediante un asistente de configuración. La estructura completa del flujo de generación y las
pantallas correspondientes a cada etapa pueden consultarse en la
Sección~\ref{subsec:mapa_navegacion}.

El funcionamiento del asistente consiste en configurar los diferentes
parámetros de generación según las necesidades del usuario. En cada etapa del
flujo se pueden modificar las opciones disponibles y avanzar al siguiente paso
mediante el botón \textit{Next}, situado en la parte inferior de la pantalla.
En caso de querer modificar una configuración anterior, es posible volver a
los pasos previos utilizando el botón \textit{Previous}.

Una vez completada la configuración, en el sexto y último paso del asistente se generan
los modelos de características solicitados. Desde esta pantalla el usuario
dispone de diferentes opciones sobre los modelos generados:
\begin{itemize}
    \item Descargar el conjunto de modelos generado mediante el botón
    \textit{Download}.
    \item Visualizar un modelo directamente desde la aplicación mediante el
    botón con el icono del ojo.
    \item Descargar individualmente cada modelo mediante el botón verde
    disponible junto a cada resultado.
    \item Abrir un modelo mediante el botón morado con el icono del
    cohete en FlamapyIDE, un entorno externo que permite visualizar, editar y
    realizar tareas de análisis automático sobre modelos de características
    (\url{https://github.com/flamapy/flamapy-ide}).
\end{itemize}

\section{Mensajes de error y resolución de incidencias}
\label{sec:errores}

Durante la configuración, el asistente valida los valores introducidos en cada
etapa antes de permitir que el usuario continúe al siguiente paso. Cuando una
opción no cumple las reglas de negocio definidas en la
Sección~\ref{sec:reglas_negocio}, se muestra un mensaje asociado al campo
correspondiente indicando la condición que debe corregirse.

Estas validaciones incluyen, entre otros casos, límites numéricos incorrectos,
distribuciones de probabilidades que no suman uno, dependencias entre niveles de
expresividad y rangos inválidos para cardinalidades, restricciones o atributos.
Una vez corregidos los valores señalados, el usuario puede continuar con el
proceso de configuración y generar los modelos solicitados.

Cuando se activa la opción \textit{Ensure satisfiable}, la generación puede
detenerse si el sistema no encuentra un modelo que el comprobador SAT confirme
como satisfacible tras 20 intentos para ese modelo. En ese caso, el asistente
muestra un mensaje de error y detiene el lote. Este resultado no demuestra que
la configuración sea insatisfacible; indica que no se ha obtenido un modelo
confirmado dentro del límite de intentos establecido.

Como indica el mensaje de la Figura~\ref{fig:error-generacion-sat}, se puede
simplificar la configuración, reducir el número de restricciones o desactivar
\textit{Ensure satisfiable} y volver a generar. Si se desactiva esta opción, el
sistema no realiza la comprobación SAT de los modelos generados.

\begin{figure}[H]
    \centering
    \includegraphics[width=1.1\textwidth]{error_generacion_sat.png}
    \caption{Mensaje mostrado cuando no se obtiene un modelo confirmado como
    satisfacible tras 20 intentos. Captura realizada por el autor.}
    \label{fig:error-generacion-sat}
\end{figure}

\section{Contacto y soporte}
\label{sec:contacto}

Para cualquier duda relacionada con el funcionamiento del sistema o para
reportar errores no contemplados en este manual, se puede contactar con el
desarrollador mediante las siguientes direcciones de correo electrónico:

\begin{itemize}
    \item antjimort@alum.us.es
    \item ajimenez34@us.es
\end{itemize}